\begin{abstract}
\noindent
This report presents \textbf{RAPID}, a Big Data cybersecurity monitoring platform designed according to the Lambda Architecture required in the project statement. The system analyzes the \textit{Cybersecurity Threat Detection Logs} dataset and combines three complementary layers: a batch layer for historical attack analysis, a speed layer for real-time detection, and a serving layer for API access and dashboard visualization.

The implemented platform uses HDFS for replayable historical storage, Apache Spark for batch and streaming computation, Apache Kafka for continuous ingestion, HBase for historical serving views, Apache Cassandra for active threats, and a Flask REST API with HTML/JavaScript dashboards for visualization. The batch layer computes malicious IP reputation, port-scan windows, attack path patterns, threat timelines, traffic-volume statistics, and multi-step attack chains. The speed layer detects brute-force attempts, known attack signatures such as SQLMap and Nikto, abnormal traffic volumes, and real-time threat scores per source IP. The serving layer merges historical HBase evidence with recent Cassandra evidence and returns decision-oriented responses with a final score and a recommendation.

The project also includes bonus analytical improvements: a Spark MLlib Random Forest classifier, multi-step attack correlation, geolocation-oriented dashboard views, and adaptive thresholds. The result is a complete academic implementation of a cybersecurity threat detection pipeline that satisfies the required batch, speed, serving, dashboard, and demonstration objectives.

\vspace{0.6em}
\noindent
\textbf{Keywords:} Big Data, Cybersecurity, Lambda Architecture, HDFS, Spark, Kafka, HBase, Cassandra, REST API, Dashboard.
\end{abstract}
